% ===== B 卷参考答案与评分标准 =====
\anstou{B}

\datiA{名词解释（每小题 5 分，共 20 分）}

\thao{1.}\da 联系是指事物内部各要素之间以及事物之间相互影响、相互作用和相互制约的关系。\pingfen{3 分}联系具有客观性、普遍性、多样性和条件性。\pingfen{2 分}

\thao{2.}\da 否定之否定规律揭示事物发展的方向和道路：事物经过两次否定、三个阶段（肯定—否定—否定之否定），\pingfen{3 分}呈现螺旋式上升或波浪式前进，方向是前进的，道路是曲折的。\pingfen{2 分}

\thao{3.}\da 真理是人们对客观事物及其规律的正确反映。\pingfen{3 分}真理具有客观性，真理与谬误有原则界限，但在一定条件下可以相互转化。\pingfen{2 分}

\thao{4.}\da 生产关系是人们在物质生产过程中形成的不以人的意志为转移的经济关系，\pingfen{2 分}包括生产资料所有制关系、生产中人与人的关系和产品分配关系。\pingfen{3 分}

\datiA{简答题（每小题 8 分，共 32 分）}

\thao{1.}\yaodian
\begin{enumerate}[label=(\arabic*),leftmargin=2.2em,itemsep=2pt,topsep=2pt]
\item 内容决定形式，有什么样的内容就要求有什么样的形式与之相适应。\pingfen{3 分}
\item 形式对内容具有反作用：适合内容的形式能促进内容发展，不适合的形式会阻碍内容发展。\pingfen{3 分}
\item 形式具有相对独立性，有时会落后于内容的变化；因此既要重视内容，也要善于选择和创造适当形式，反对形式主义和形式虚无主义。\pingfen{2 分}
\end{enumerate}

\thao{2.}\yaodian
\begin{enumerate}[label=(\arabic*),leftmargin=2.2em,itemsep=2pt,topsep=2pt]
\item \textbf{区别}：感性认识是对事物现象、外部联系的反映，形式为感觉、知觉、表象；理性认识是对事物本质、内部联系的反映，形式为概念、判断、推理。\pingfen{4 分}
\item \textbf{相互过渡}：感性认识有待于上升为理性认识，这需要掌握丰富而真实的感性材料并运用科学思维方法进行加工；理性认识又必须回到实践中接受检验并指导实践。\pingfen{3 分}
\item 二者在实践基础上统一，割裂会走向唯理论或经验论。\pingfen{1 分}
\end{enumerate}

\thao{3.}\yaodian
\begin{enumerate}[label=(\arabic*),leftmargin=2.2em,itemsep=2pt,topsep=2pt]
\item 规律内容：生产力决定生产关系，决定其性质和变革方向；生产关系对生产力具有反作用，适合时促进、不适合时阻碍。\pingfen{4 分}
\item 这一矛盾运动是社会发展的基本动力，生产关系一定要适合生产力状况是社会发展的普遍规律。\pingfen{2 分}
\item 现实意义：判断一种生产关系是否先进，根本标准是看它是否适合生产力发展；改革就是调整生产关系中不适合生产力的部分。\pingfen{2 分}
\end{enumerate}

\thao{4.}\yaodian
\begin{enumerate}[label=(\arabic*),leftmargin=2.2em,itemsep=2pt,topsep=2pt]
\item 经济基础决定上层建筑，上层建筑反作用于经济基础；上层建筑一定要适合经济基础状况是社会发展的客观规律。\pingfen{3 分}
\item 适合时，上层建筑为经济基础服务，促进生产力发展和社会进步。\pingfen{2 分}
\item 不适合时，上层建筑会成为经济基础变革的阻碍，此时就要通过革命或改革调整上层建筑。\pingfen{2 分}
\item 上层建筑具有相对独立性，可能超前或滞后于经济基础的变化，这是需要主动改革的原因。\pingfen{1 分}
\end{enumerate}

\datiA{论述题（每小题 15 分，共 30 分）}

\thao{1.}\yaodian
\begin{enumerate}[label=(\arabic*),leftmargin=2.2em,itemsep=2pt,topsep=2pt]
\item \textbf{原理}：矛盾的普遍性即矛盾存在于一切事物中并贯穿始终；矛盾的特殊性即具体事物的矛盾各有其特点。\pingfen{4 分}
\item \textbf{关系}：普遍性寓于特殊性之中，特殊性中包含普遍性，二者在一定条件下相互转化。\pingfen{4 分}
\item \textbf{方法论}：坚持共性与个性、一般与个别的统一，反对只讲共性不讲个性的教条主义，也反对只讲个性否认共性的经验主义。\pingfen{3 分}
\item \textbf{联系中国特色社会主义}：社会主义的一般原则（共性）必须同中国具体国情（个性）相结合，照搬别国模式行不通；中国特色社会主义正是共性寓于个性的产物，这也是“中国化”的哲学依据。\pingfen{4 分}
\end{enumerate}

\thao{2.}\yaodian
\begin{enumerate}[label=(\arabic*),leftmargin=2.2em,itemsep=2pt,topsep=2pt]
\item \textbf{确定性}：实践是检验真理的唯一标准，此外别无标准；凡经实践证实的认识都具有客观真理性。\pingfen{4 分}
\item \textbf{不确定性}：一定历史阶段的实践有其局限性，只能在一定程度上检验认识；实践检验往往需要一个过程，不是一次完成。\pingfen{4 分}
\item \textbf{二者统一}：确定性是绝对的、无条件的，不确定性是相对的、有条件的；承认不确定性不是否定实践标准，而是要求在实践中不断检验和发展认识。\pingfen{4 分}
\item \textbf{意义}：既反对怀疑实践标准的相对主义，也反对把某一阶段的实践结论绝对化的僵化态度。\pingfen{3 分}
\end{enumerate}

\datiA{材料分析题（共 18 分）}

\thao{1.}
\textbf{（1）}\pingfen{9 分}
\begin{enumerate}[label=(\arabic*),leftmargin=2.2em,itemsep=2pt,topsep=2pt]
\item 认识来源于实践。照搬文献路线失败，说明他人实践得出的认识不能直接代替自己的实践检验。\pingfen{3 分}
\item 认识过程是“实践—认识—再实践—再认识”的循环往复、不断上升的过程。团队通过三百余组实验，由现象到本质，逐步形成对工艺规律的正确认识。\pingfen{3 分}
\item 实践是检验认识真理性的唯一标准。“能不能用，还得我们自己试”正是这一原理的朴素表达；认识最终要回到实践中指导实践。\pingfen{3 分}
\end{enumerate}

\textbf{（2）}\pingfen{9 分}
\begin{enumerate}[label=(\arabic*),leftmargin=2.2em,itemsep=2pt,topsep=2pt]
\item 真理尺度指认识必须符合客观事物及其规律，即“求真”；价值尺度指实践必须满足人的需要，即“求善”。\pingfen{3 分}
\item 成功的实践是二者的统一：只符合规律而不满足需要，成果无意义；只讲需要而不顾规律，必然在实践中失败。\pingfen{3 分}
\item 科技创新中，性能达标体现真理尺度，成本可控、可产业化、安全环保体现价值尺度；二者的统一是科技成果转化为现实生产力的前提。\pingfen{3 分}
\end{enumerate}
